\documentclass[../main.tex]{subfiles}

\begin{document}
% -----------------------------
\section{Simulator}
% -----------------------------
A simulator is a component that models the behavior of a system by executing its operations, sometimes including
hardware-level details such as memory, timing, or architectural constraints. In this tutorial, we will use an interpreter,
which directly executes our IR operations to validate correctness without modeling low-level hardware behavior.

Create 'simulator' folder next to your \path{hc_main.py} and \path{hc_dialect.py}:
\begin{lstlisting} [language=bash, caption={Simulator folder structure}]
simulator
├── __init__.py
├── interpreter.py
\end{lstlisting}
First add helper functions:
\begin{lstlisting} [language=Python, caption={Helper functions}]
from typing import Any
from xdsl.ir import Operation, SSAValue
from xdsl.dialects import func

# -----------------------------
# Helpers to read constants
# -----------------------------

def _int_from_attr(attr: Any) -> int | None:
    """
    Extract an int from common xdsl attribute shapes (IntegerAttr/IntAttr, etc.).
    """
    if attr is None:
        return None

    # Often: attr.value is IntAttr; or attr.data is int; or nested.
    for k in ("data", "value"):
        if hasattr(attr, k):
            inner = getattr(attr, k)
            # inner might again be an attr
            for k2 in ("data", "value"):
                if hasattr(inner, k2):
                    try:
                        return int(getattr(inner, k2))
                    except Exception:
                        pass
            try:
                return int(inner)
            except Exception:
                pass

    try:
        return int(attr)
    except Exception:
        return None


def _const_value(op: Operation) -> int | None:
    """
    arith.constant -> python int
    """
    if op.name != "arith.constant":
        return None
    v = getattr(op, "value", None)
    return _int_from_attr(v)
\end{lstlisting}

Create Interpreter class that will let us instantiate our interpreter:
\begin{lstlisting} [language=Python, caption={Interpreter class}]
# -----------------------------
# Interpreter
# -----------------------------

class Interpreter:
    def __init__(self):
        # SSAValue is not always hashable across versions -> use id()
        self.env: dict[int, int] = {}

    def _get(self, v: SSAValue) -> int:
        return self.env[id(v)]

    def _set(self, v: SSAValue, value: int) -> None:
        print(f"Set SSAValue: {v} = {int(value)}")
        self.env[id(v)] = int(value)

\end{lstlisting}
Implement \path{run_block} method inside of this class. It deals with both, \path{arith} and \path{hc} instructions. 
\begin{lstlisting} [language=Python, caption={run\_block method implementation}]
def run_block(self, block) -> int:
    """
    Execute a single basic block (straight-line).
    Returns the integer from func.return.
    """
    for op in list(block.ops):
        name = op.name

        # --- constants ---
        if name == "arith.constant":
            val = _const_value(op)
            if val is None:
                raise RuntimeError(f"Couldn't read constant value from op: {op}")
            self._set(op.results[0], val)
            continue

        # --- high-level HC ops ---
        if name == "hc.add":
            a, b = op.operands
            self._set(op.results[0], self._get(a) + self._get(b))
            continue

        if name == "hc.sub":
            a, b = op.operands
            self._set(op.results[0], self._get(a) - self._get(b))
            continue

        if name == "hc.mul":
            a, b = op.operands
            self._set(op.results[0], self._get(a) * self._get(b))
            continue

        if name == "hc.relu":
            (x,) = op.operands
            self._set(op.results[0], max(self._get(x), 0))
            continue

        # --- lowered arith ops (so you can run after lowering too) ---
        if name == "arith.addi":
            a, b = op.operands
            self._set(op.results[0], self._get(a) + self._get(b))
            continue

        if name == "arith.subi":
            a, b = op.operands
            self._set(op.results[0], self._get(a) - self._get(b))
            continue

        if name == "arith.muli":
            a, b = op.operands
            self._set(op.results[0], self._get(a) * self._get(b))
            continue

        if name == "arith.maxsi":
            a, b = op.operands
            self._set(op.results[0], max(self._get(a), self._get(b)))
            continue

        # --- return ---
        if name == "func.return":
            if len(op.operands) != 1:
                raise RuntimeError(f"Expected func.return with 1 operand, got: {op}")
            return self._get(op.operands[0])

        raise RuntimeError(f"Unsupported op in interpreter: {name}\nOp: {op}")

    raise RuntimeError("Block ended without func.return")
\end{lstlisting}
This method executes a single basic block in a straight-line manner, meaning it processes operations sequentially
without handling control flow such as branches or loops. It iterates through each operation in
the block and dispatches behavior based on the operation’s name.

For constants (arith.constant), it extracts the literal value and stores it in the
interpreter’s internal environment. The interpreter maintains a mapping from SSA values
to concrete integer values.

For high-level operations (hc.add, hc.sub, hc.mul, hc.relu), it retrieves operand
values, performs the corresponding
Python arithmetic, and stores the computed result back into the environment.
This directly implements the semantics of our high-level dialect.

It also supports lowered arithmetic operations
(
\texttt{arith.addi},
\texttt{arith.subi}, \texttt{arith.muli}, \texttt{arith.maxsi}),
which allows the same interpreter to run code after lowering passes. This makes the interpreter
useful both before and after middle-end transformations, reinforcing
its role as a validation tool.

Next, implement \path{run_module} that calls \path{run_block}:
\begin{lstlisting} [language=Python, caption={run\_module method implementation}]
def run_module(self, module, func_name: str = "my_func") -> int:
    """
    Find func.func @func_name and execute its first block.
    """
    # module.body.blocks[0].ops usually contains top-level ops
    for top_block in module.body.blocks:
        for op in list(top_block.ops):
            if isinstance(op, func.FuncOp):
                # Try to read the symbol name
                sym = getattr(op, "sym_name", None)
                name = None
                if sym is not None:
                    name = _int_from_attr(getattr(sym, "data", None))  # probably not int, ignore
                    # better: sym.data is usually a string
                    if hasattr(sym, "data"):
                        name = sym.data
                    else:
                        name = str(sym)

                if name == func_name or (name is None and func_name == "my_func"):
                    # assume single-block body for now
                    body_block = list(op.body.blocks)[0]
                    return self.run_block(body_block)

    raise RuntimeError(f"Function not found: {func_name}")
\end{lstlisting}
This method searches the given module for a function with the specified name
and executes it using the interpreter. It iterates over all top-level blocks
in the module and then over their operations, looking specifically for
instances of func.FuncOp, which represent function definitions in the IR.

For each function operation found, it attempts to extract the
function’s symbolic name (\path{sym_name}). Since attributes in xDSL are
often wrapped objects, the code carefully accesses sym.data
(which typically holds the actual string name) and falls back to
string conversion if necessary.

Once the function name matches \path{func_name} (defaulting to \path{my_func}),
the method assumes the function body contains a single block, retrieves that
first block, and passes it to \path{self.run_block()} for execution.

To use it in our \path{hc_main.py}:
\begin{lstlisting} [language=bash, caption={Instantiating Interpreter}]
from simulator.interpreter import Interpreter
...
def main():
    # <Building module>
    ...

    hc_interpreter = Interpreter()
    result_before = hc_interpreter.run_module(module, "my_func")
    print("\nResult (interpreted, before lowering):", result_before)

    # <Middle-end lowering part>
    ...

    lower_interpreter = Interpreter()
    result_after = lower_interpreter.run_module(module, "my_func")
    print("\nResult (interpreted, after lowering):", result_after)
    \end{lstlisting}

\end{document}
